% Independent full proofs for the captured original C+ formal class.
\section{C+ laws and finite histories: proofs}
\label{sec:annex-p16-cplus-action-language}
Use \cref{sec:p16-cplus-action-language} and the activated-head semantics
of \cref{def:p16-nct-theory}. The source counterparts are Propositions
7--9 and their Appendix proofs in the retained author manuscript
\cite{GiunchigliaEtAl2003CplusAuthor}.
Formula heads are not assumed atomic.

\begin{annexitem}{annex:p16.cplus-paths}{thm:p16-cplus-paths}{Coordinate decomposition and endpoint states}
\begin{proof}
Fix a candidate history \(I\) and activate all rule bodies in \(D_m\).
Let \(\Delta_j\) be the static heads activated at state coordinate \(j\),
\(\Pi_0\) the actual-value atoms of every simple fluent at time zero,
\(A_i\) the action dynamic heads activated at event coordinate \(i\),
and \(B_j\), \(1\leq j\leq m\), the fluent dynamic heads activated
for successor \(j\). Stamps are implicit in these block names.
The reduct is the union of
\[
 \Delta_0\cup\Pi_0,\qquad
 A_i\ (0\leq i<m),\qquad
 \Delta_j\cup B_j\ (1\leq j\leq m).
\]
Each block contains formulas on only its own coordinates.
Every \(B_j\) contains only simple fluents; every \(A_i\) contains
only action constants. Mixed bodies have already been evaluated, and
introduce no remaining cross-coordinate head dependencies.

For these disjoint coordinate blocks, \(I\) is the unique global model
iff its restriction is the unique model of each block.
If it is unique globally, each restriction satisfies its block;
another block model, combined with the unchanged other restrictions,
would contradict global uniqueness. Conversely, each block's unique
model fixes every coordinate, and their union satisfies the reduct.
This argument permits arbitrary Boolean formulas, empty blocks and
singleton domains, and does not assume a finite signature.

We need one additional observation. Suppose \(s\) uniquely satisfies
\(\Delta\cup B\), where \(B\) mentions only simple fluents, and
\(\Delta\) is the static-head set activated at \(s\).
Let \(\Pi(s)\) fix all its simple fluent values.
Then \(s\) uniquely satisfies \(\Delta\cup\Pi(s)\).
It satisfies this set; any other model agrees on every simple fluent,
hence satisfies all formulas of \(B\) because \(s\) does.
It would be another model of \(\Delta\cup B\), a contradiction.
Thus \(s\) is a state even though its simple values in the history may
have been explained by dynamic laws rather than initial self-rules.

For \(m=1\), decomposition says the origin satisfies its state block
uniquely, and the observation makes the target a state.
This proves the endpoint assertion.

For general \(m>0\), a global explanatory history has a state at
coordinate zero and, by the observation, at every later coordinate.
For each adjacent triple, its local \(D_1\) reduct has origin block
\(\Delta_i\cup\Pi(s_i)\), its event block \(A_i\), and successor block
\(\Delta_{i+1}\cup B_{i+1}\). All three are unique by the preceding
arguments, so every triple is a transition.
Conversely, a path's local transitions make every event block and
every successor block unique. Its initial state supplies the unique
first block. Therefore their global union uniquely models \(D_m\).
For \(m=0\) the assertion is exactly the definition of a state.
\end{proof}
\end{annexitem}

\begin{annexitem}{annex:p16.cplus-head-exception}{thm:p16-cplus-head-exception}{An invalid derived dynamic head}
\begin{proof}
There are no simple fluents, hence no initial self-rule.
The zero-horizon theory is \(0:p\Leftarrow0:p\).
At true \(p\) the reduct is \(\{0:p\}\), uniquely true; at false
\(p\) it is empty and has both Boolean assignments.
The proposed one-step theory is
\[
 0:p\Leftarrow0:p,\quad1:p\Leftarrow1:p,\quad
                                  1:\neg p\Leftarrow0:p.
\]
At \(p_0=t,p_1=f\) its reduct is \(\{0:p,1:\neg p\}\),
which fixes both coordinates uniquely.
The false target is not a zero-horizon state. The failed step in the
general proof is the assertion that dynamic heads mention only simple
coordinates: fixing simple values now does not preserve those heads.
\end{proof}
\end{annexitem}

\begin{annexitem}{annex:p16.cplus-embedding}{thm:p16-cplus-embedding}{The static embedding}
\begin{proof}
At horizon zero there are no initial self-rules because all fluents
are derived. Removing the zero stamp identifies \(D_0\) with \(T\),
including body activation and unique head satisfaction.
Thus state assignments and explanatory models coincide.
At horizon one there are exactly two disjoint copies of \(T\), with
no action coordinates and no dynamic rules. A pair of assignments
uniquely satisfies its activated heads iff each assignment does so
in its own copy. This gives every ordered pair of states and no others.
\end{proof}
\end{annexitem}

\begin{annexitem}{annex:p16.cplus-rigid}{thm:p16-cplus-rigid}{Compression and both failures}
\begin{proof}
Under the hypotheses, a head mentioning \(R\) can only be static:
action heads have no fluents, and fluent dynamic heads have no
derived fluents. Such a static law contains only \(R\) throughout.
Its activated heads form a pure-\(R\) block at each state time.
All other activated heads mention only non-\(R\) constants, though
their evaluated bodies may have tested \(R\). There are no initial
simple self-rules on \(R\).

If \(I\) is a model, the path theorem and rigidity give equal values
of every \(R\)-constant at every time. Thus activation of pure-\(R\)
laws gives identical copies of a formula theory \(\Delta_R\).
All remaining activated heads form a theory \(\Gamma\) on the other
coordinates. The old reduct is disjoint copies of \(\Delta_R\), one
per state time, together with \(\Gamma\).
Its uniqueness is equivalent to uniqueness of each copy and of
\(\Gamma\). Erasing stamps gives exactly the reduct at \(I^R\):
one \(\Delta_R\) block and the same \(\Gamma\), with the same
activation of every body. Hence \(I^R\) is explanatory.

Conversely, let \(J\) be explanatory for the compressed theory.
Lift its shared \(R\)-values to every state time, obtaining \(I\).
Body activation agrees under substitution. The compressed reduct is
one pure-\(R\) block and a non-\(R\) block, uniquely satisfied.
Replacing its first block by identical disjoint copies preserves
uniqueness and satisfaction, so the lift is explanatory for \(D_m\).
Lifting and compression are inverse. This includes \(R=\varnothing\)
and \(m=0\).

For the first counterexample, the simple Boolean \(p\) has static
\(p\Leftarrow p\) at each time, initial self-rules for both values,
and rigidity constraints forbidding unequal consecutive values.
At horizon one unequal values activate a false head.
With both values true, the reduct fixes both true.
With both false, only the initial negative self-rule is active;
the successor has no explanation and is not uniquely fixed.
After merging, constraints are never active and both initial
self-rules are available at the one shared coordinate.
The false candidate has reduct \(\{\neg p\}\), uniquely false.
The first hypothesis has therefore prevented a real spurious model.

For the second example the original one-step theory is
\[
\begin{gathered}
 p_0\Leftarrow q_0,\quad p_1\Leftarrow q_1,\\
 q_0\Leftarrow q_0,\quad\neg q_0\Leftarrow\neg q_0,\quad
 q_1\Leftarrow q_1,\quad\neg q_1\Leftarrow\neg q_1.
\end{gathered}
\]
A state with \(q=f\) leaves \(p\) unexplained, while \(q=t\)
forces and explains \(p=t\). Hence there is only the true/true state
and \(p\) is rigid in the actual transition system.
After merging \(p_0,p_1\), the candidate \(p=t,q_0=f,q_1=t\)
has reduct \(\{p,\neg q_0,q_1\}\), uniquely satisfied.
Its unmerged lift leaves \(p_0\) unexplained and is not explanatory.
The head law for \(p\) had depended on \(q\notin R\), violating the
second hypothesis.
\end{proof}
\end{annexitem}

\begin{annexitem}{annex:p16.cplus-outcomes}{thm:p16-cplus-outcomes}{Source case derivations}
\begin{proof}
All descriptions in this item are definite after abbreviation expansion;
use exact completion or, equivalently, inspect their activated value
atoms. Initial simple-value self-rules fix the initial assignment;
action exogeneity fixes the chosen event assignment. It remains to
show satisfaction and uniqueness on successor coordinates.

For the basic action, positive successor support is
\((p'\land p)\lor a\), negative support \(\neg p'\land\neg p\).
If \(a=t\), positive support forces \(p'=t\); if \(a=f\), changing
the old value has no support, while its unchanged value has inertia.
Both unchanged and forced-true reducts uniquely fix the successor.
Thus the four input pairs give exactly \(p'=p\lor a\).

For the specified monkey path, take these successive full fluent
assignments:
\[
\begin{array}{c|ccc|cc}
 &Loc(Monkey)&Loc(Box)&Loc(Bananas)&OnBox&HasBananas\\ \hline
 s_0&L_1&L_1&L_3&f&f\\
 s_1&L_3&L_3&L_3&f&f\\
 s_2&L_3&L_3&L_3&t&f\\
 s_3&L_3&L_3&L_3&t&t.
\end{array}
\]
At \(s_0\) no static antecedent requiring OnBox or HasBananas holds,
and initial self-rules fix the displayed values.
Push to \(L_3\) is allowed at \(s_0\): the monkey is at the box,
not at the destination and not on it. Its two direct causes fix the
changed locations, inertia fixes all other values.
ClimbOn is allowed at \(s_1\), fixes OnBox true, and the activated
static monkey-location law agrees with the location's inertia.
Grasp is allowed at \(s_2\), fixes HasBananas true, and both static
location schemas agree with the unchanged locations.
All listed concurrency prohibitions have false antecedents because
only the indicated action is true. Every activated head is a value
atom agreeing with the displayed successor, and every coordinate
has such an active atom from a direct effect or inertia.
The reduct therefore has exactly that successor.
The path theorem joins these three transitions.

For shooting, the old target is \(Turkey_1\), so the active direct heads
are \(\neg Loaded'\) and \(\neg Alive(Turkey_1)'\).
The other turkey and target are fixed by inertia.
The shooting prohibition is inactive because Loaded was true;
all aim/shoot conjunctions are false. The displayed successor
satisfies all active atoms and is uniquely fixed.
Attempted survival of \(Turkey_1\) or retention of Loaded conflicts with
its respective direct negative head.

For travel, Go(Work) forces Jack to Work.
The old equality of car and Jack activates the eligibility condition
of the may-cause law. If car stays Home, Home inertia is active and
the optional Work self-test is false. If car moves to Work, its
optional Work explanation is active and Home inertia is false.
Both reducts uniquely fix their respective coordinates.
The car has only these two values; Jack has only the forced Work
value, proving exactly two successors, without weighting them.

For the table, a selected lift forces its end High; an unselected
end can only remain Low by inertia. If the candidate left an unlifted
end High, it would have no active value head and fail uniqueness.
With one lift the unequal levels cause \(\neg OnTable'\);
a true candidate would also activate positive inertia and be
inconsistent. With both lifts the equal levels supply no negative
static head. True OnTable has inertia; false OnTable has no support
and fails uniqueness. Thus each of these event choices has precisely
the claimed successor.

For the pendulum suppose old Right is true. Without Hold, the false
successor activates the negative flip default and is uniquely fixed;
the true successor activates no head and fails uniqueness.
With Hold, the positive effect is active, making a false successor
inconsistent with its activated flip head and explaining the true
successor uniquely. The case of old false is obtained by explicitly
interchanging positive and negative heads: without Hold only the
positive flip default explains a successor, and with Hold the
negative effect forces and explains persistence.
\end{proof}
\end{annexitem}

\begin{annexitem}{annex:p16.cplus-definitions}{thm:p16-cplus-definitions}{Fresh definitions and attribute counts}
\begin{proof}
For fresh derived \(d\) there is no initial self-rule.
At each state time the added laws supply \(d\Leftarrow F\) and
\(\neg d\Leftarrow\neg d\). If \(F\) is true and \(d\) false,
both heads activate and are inconsistent.
If \(F\) is false and \(d\) true, neither added head activates.
Varying that coordinate alone then preserves all activated old
heads, since they contain no \(d\), contradicting uniqueness.
Consequently every explanatory model has \(d\equiv F\).
At that assignment exactly one added value head is active, fixing
\(d\), and the other activated heads are the old reduct.
This gives both directions and uniqueness of extension for states
and finite histories. With replacements in old conditions, activation
of each old head agrees because \(d\) and \(F\) have identical truth
values; old heads still do not contain \(d\). The same reasoning
therefore applies. This argument does not require the old heads to
be atomic.

For failure with a simple \(d\), take a simple Boolean old \(p\),
let \(F=p\), and take no old laws. At horizon zero the assignment
\(p=f,d=t\) activates the initial self-rules \(\neg p,d\), but
neither defining head. Its reduct uniquely fixes this assignment,
despite \(d\not\equiv F\).

For publication events the two always-laws exclude precisely violations
of \(Kind=None\) iff Publish is false and \(Length=None\) iff Publish
is false. With false Publish only the single pair None,None remains.
With true Publish, Kind has two choices and Length one hundred;
their Cartesian product has \(200\) assignments.
Each is explained by action exogeneity.
From the all-false state, Publish supplies HasPublications.
Journal supplies HasJournalPublications, and any numbered length
greater than 30 supplies HasLongPublications.
A fluent with a positive direct cause must be true; otherwise its
only supported value is the old false value from inertia.
Thus Journal/31 gives three true values, while Conference/20 gives
true,false,false, and each respective reduct uniquely fixes them.
The event constraints are satisfied by both choices.
\end{proof}
\end{annexitem}

